\section{人生哲学：长期主义、初心与勇气}

\subsection{长期主义的数学基础}

陈天奇在《机器学习科研这十年》中写道：``很多困难和无助只是随机涨落的一部分，只要付出足够多的时间和耐心，随机过程总会收敛到和付出相称的稳态。''

他说这句话``就是一个定理''------随机过程的 time average 确实会收敛到稳态分布。但知道定理是一回事，在挫折中坚持是另一回事。

\begin{importantbox}{接受失败的能力}
``失败也没那么可怕，然后要接受可以失败这个事实。''这个信念从高中 OI 只拿二等奖、到本科做深度学习两年无果、到创业的各种转型，一直在被强化。陈天奇认为：\textbf{经历越多失败，就越敢放手做事}------因为``失败了也没什么大不了的，我还是活得好好的''。早期失败甚至是一种``保护''：如果一开始就太成功，反而可能形成路径依赖。
\end{importantbox}

\subsection{勇气来自过去的自己}

陈天奇的一个深刻洞察是：\textbf{做事情的勇气并不随经验增长而自动增加}。恰恰相反，看到的东西越多，反而越容易畏首畏脚。

``我觉得很多时候你的勇气可能是你过去的自己给你自己的约定。你要达成过去自己对自己的约定，才要继续去坚持做这个 style 的工作。''

他把这个过程描述为一个循环：过去的自己给当下的自己勇气，当下的自己在做的过程中不断鼓励自己，再把这份勇气交给未来的自己。

\subsection{保持初心是最难的事}

当被问到``展望未来 5--10 年想干什么''时，陈天奇的回答出人意料：

\textbf{``我觉得现在的我最难的事情，还是要保持初心。''}

他解释说，随着资源变多、选择变多、责任变大，反而更难保持当年``出生牛犊不怕虎''的心态。做 MXNet 时完全不担心 TensorFlow 的竞争，但现在面对一个新项目时会不由自主地问``它到底能不能成功''。保持初心需要刻意的自我提醒。

主持人追问：你能调动的资源比 10 年前多了很多个数量级，为什么反而觉得保持初心更难？陈天奇回答：``拿到资源的同时必然代表着责任。你必须保证团队可以成功，必须产出相应的 impact 来 justify 做相关的事情。如果被这些压力裹挟而走，就会失去做 risky 事情的勇气。''

他的应对策略是：在 CMU 团队中继续做更 risky 的事情，``愿意放弃和愿意接受失败------这个心态其实和一开始是一样的''。资源和团队变大了，能做的事的确比原来更多了，但心态上的转变可能让人``忘记自己的那个心态''。

\begin{warningbox}{成功后的心态陷阱}
这段讨论揭示了一个深刻的悖论：成功和经验的积累可能\textbf{侵蚀做创新所需的勇气}。看到的东西越多、知道的风险越多，反而越容易保守。陈天奇认为，对抗这种倾向的方法是回到初心------不是用理性的投入产出分析来做决策，而是问自己：``这是不是我理想中想要走的路？''如果是，就不管结果如何都去做。
\end{warningbox}

\subsection{关于成功与幸福}

主持人在最后部分问了几个个人化的问题。

\textbf{关于成功}：陈天奇认为成功有很多定义------影响力、财富、技术突破等。但对他而言，最重要的是``做问心无愧的事情''。``能够在技术上面可以影响大家，并且可以为这个世界一起做一些贡献，就是一件好事情。''

\textbf{关于对未来自己的寄语}：如果让他给过去和未来的自己各说一句话，他的选择不是回头改变什么，而是相反------\textbf{他需要过去的自己来提醒现在和未来的自己}：记住对自己的承诺，坚持理想，往下走下去。

\textbf{关于理想的人生状态}：10 年后 20 年后理想的陈天奇，应该是``有更多次失败的一个陈天奇------并且可以继续保持那一个最初的心态去挑战我们想要挑战的东西''。

\subsection{人生哲学本章小结}

陈天奇的人生哲学可以用三个关键词概括：\textbf{长期主义}（相信随机过程会收敛）、\textbf{初心}（做问心无愧的事）、\textbf{勇气}（接受失败，不被资源和责任裹挟）。这些不是空洞的口号，而是在 20 年的科研和创业中被反复验证的行动准则。值得注意的是，他并不认为这些品质随着经验增长会自然增强------恰恰相反，他认为保持初心和勇气需要持续的、刻意的自我提醒。

